\begin{textsample}{NEG Dim 1 – global\_south\_2024\_10 – Score 1.00 – global\_south\_2024\_10/115686221.txt}  \label{ex:f1_neg_436}
Israel ' s bombing of Gaza is exposing millions of people to the long-term threat of asbestos, a hazardous substance released as buildings are destroyed.
Once airborne, asbestos fibres can be inhaled, causing cancer and other deadly diseases.
Experts fear the health consequences for the people of Gaza will last for decades.
The United Nations estimates that around 800,000 tonnes of debris across the Gaza Strip may be contaminated with asbestos, making the situation even more dire for its 2.3 million residents.
Asbestos is usually harmless when undisturbed, but becomes highly carcinogenic when released into the air, posing significant health risks.
Roger Willey, a leading expert on asbestos, described the situation as a"death sentence"for the people of Gaza.
He explained :"The best thing to do if asbestos becomes airborne is to drive as far away as possible.
But for the people of Gaza, that ' s simply impossible."Israel ' s bombing campaign has resulted in an estimated 42 million tonnes of debris across Gaza since October 2023, and much of it @ @ @ @ @ @ @ @ @ @ to the UN Environment Programme.
The destruction has left an already vulnerable population exposed to these long-term health hazards.
Liz Darlison, CEO of Mesothelioma UK, a charity supporting those affected by asbestos-related diseases, said,"The long-term effects of asbestos exposure will constitute a tragedy that will unfold in the years ahead.
There is no safe level of exposure."She added that the legacy of this war will continue for many years, with cases of cancer likely to be reported"for decades."Drawing parallels to the aftermath of the 9/11 attacks, Willey pointed out that asbestos exposure around the World Trade Center has caused more than 4,000 deaths due to related illnesses.
He warned that Gaza faces a similar fate.
The World Health Organization ( WHO ) has also raised concerns, stating that dust from the bombed-out buildings is spreading hazardous materials, posing serious health risks.
With only 11 \% of Gaza remaining a"safe zone"according to UN estimates, residents have little means of escaping @ @ @ @ @ @ @ @ @ @, the International Committee of the Red Cross ( ICRC ) has warned of the pervasive risk of unexploded ordnance across Gaza.
The scale of destruction has left more than 128,000 buildings destroyed or severely damaged, exacerbating the risk of asbestos exposure.
Gaza, already besieged and suffering, now faces the silent, deadly threat of asbestos, compounding the long-term impact of Israel ' s ongoing military campaign.

% matched lemmas: 
\end{textsample}
